\begin{proof}[Proof of \cref{obj:lem:finite-polynomial-duality}]
\begin{enumerate}
\item For a continuous function \(f:[v_-,v_+]\to\mathbb R\), write
\[
\|f\|_\infty=\sup_{v\in[v_-,v_+]}|f(v)|,
\qquad
\mathscr P_L=\{Q:\ Q\text{ is a real polynomial of degree at most }L\}.
\]
The restrictions of the polynomials in \(\mathscr P_L\) form a finite-dimensional normed space. Hence the ball
\[
\mathscr B_{\mathrm{app}}
=
\{Q\in\mathscr P_L:\ \|Q\|_\infty\le 2\|g\|_\infty+1\}
\]
is compact, and the continuous function \(Q\mapsto\|g-Q\|_\infty\) attains its minimum there. Choose a minimizer \(Q_*\). Since the zero polynomial belongs to this ball,
\[
\|g-Q_*\|_\infty\le \|g\|_\infty.
\]
For every \(Q\in\mathscr P_L\setminus\mathscr B_{\mathrm{app}}\), the triangle inequality gives
\[
\|g-Q\|_\infty
\ge \|Q\|_\infty-\|g\|_\infty
> \|g\|_\infty+1
> \|g-Q_*\|_\infty.
\]
Thus \(Q_*\) is a global minimizer. Define the approximation error and residual by
\[
\mathcal E_{\mathrm{app}}
=
\inf_{Q\in\mathscr P_L}\|g-Q\|_\infty
=
\|g-Q_*\|_\infty\ge0,
\qquad
\rho(v)=g(v)-Q_*(v)
\quad(v\in[v_-,v_+]).
\]
In particular, \(\rho\) is continuous and
\[
|\rho(v)|\le\mathcal E_{\mathrm{app}}
\quad(v\in[v_-,v_+]).
\]
% lean: Causalean.Mathlib.Analysis.FinitePolynomialAlternationDuality.exists_bestPolynomial

\item If \(\mathcal E_{\mathrm{app}}=0\), define
\[
v_i=v_-+\frac{i-1}{L+1}(v_+-v_-)
\quad(1\le i\le L+2),
\qquad
\eta=1.
\]
The denominator is positive for every \(L\ge0\), so these nodes are strictly increasing and belong to \([v_-,v_+]\). Moreover, the residual bound implies \(\rho(v)=0\) throughout the interval. Consequently,
\[
\rho(v_i)=\eta(-1)^{i-1}\mathcal E_{\mathrm{app}}
\quad(1\le i\le L+2).
\]
We next obtain the same identity when \(\mathcal E_{\mathrm{app}}>0\).
% lean: Causalean.Mathlib.Analysis.FinitePolynomialAlternationDuality.exists_equioscillationWitness

\item Suppose \(\mathcal E_{\mathrm{app}}>0\), and assume for contradiction that there are no \(L+2\) strictly increasing interval points at which \(\rho\) alternates between its two extreme values. Define
\[
\mathcal K_{\mathrm{ext}}
=
\{v\in[v_-,v_+]:|\rho(v)|=\mathcal E_{\mathrm{app}}\},
\qquad
\mathcal K_\pm
=
\{v\in\mathcal K_{\mathrm{ext}}:
\rho(v)=\pm\mathcal E_{\mathrm{app}}\}.
\]
These sets are compact; \(\mathcal K_{\mathrm{ext}}\) is nonempty because \(|\rho|\) attains its maximum, and \(\mathcal K_+\) and \(\mathcal K_-\) are disjoint.

We construct a polynomial \(Q_{\mathrm{sgn}}\in\mathscr P_L\) satisfying
\[
\rho(v)Q_{\mathrm{sgn}}(v)>0
\quad(v\in\mathcal K_{\mathrm{ext}}).
\]
If \(\mathcal K_+=\varnothing\), take \(Q_{\mathrm{sgn}}(v)=-1\); then the product equals \(\mathcal E_{\mathrm{app}}>0\) on \(\mathcal K_{\mathrm{ext}}\). If \(\mathcal K_+\ne\varnothing\) and \(\mathcal K_-=\varnothing\), take \(Q_{\mathrm{sgn}}(v)=1\), with the same conclusion.

It remains to consider the case in which both sign sets are nonempty. Their compactness and disjointness provide a number \(\delta_{\mathrm{ext}}>0\) such that
\[
|v-v'|>\delta_{\mathrm{ext}}
\quad(v\in\mathcal K_+,\ v'\in\mathcal K_-).
\]
Set
\[
\varepsilon_{\mathrm{ext}}=\frac{\delta_{\mathrm{ext}}}{5}.
\]
Choose a finite subcover of \(\mathcal K_{\mathrm{ext}}\) by open balls of radius \(\varepsilon_{\mathrm{ext}}\) centered in \(\mathcal K_{\mathrm{ext}}\), and list the distinct centers in increasing order:
\[
c_1<\cdots<c_{N_{\mathrm{cov}}},
\qquad
c_j\in\mathcal K_{\mathrm{ext}},
\qquad
\mathcal K_{\mathrm{ext}}
\subseteq
\bigcup_{j=1}^{N_{\mathrm{cov}}}
(c_j-\varepsilon_{\mathrm{ext}},c_j+\varepsilon_{\mathrm{ext}}),
\qquad
N_{\mathrm{cov}}\ge1.
\]
Define their signs and the number of adjacent sign changes by
\[
s_j=\frac{\rho(c_j)}{\mathcal E_{\mathrm{app}}}\in\{-1,1\},
\qquad
d_{\mathrm{flip}}
=
\#\{j:\ 1\le j<N_{\mathrm{cov}},\ s_j\ne s_{j+1}\}.
\]
If \(d_{\mathrm{flip}}>L\), selecting one center from each of the first \(L+2\) successive constant-sign runs gives \(L+2\) increasing alternating extrema, contrary to the supposition. Therefore
\[
d_{\mathrm{flip}}\le L.
\]

Define polynomials backwards along the list of centers:
\begingroup\footnotesize
\[
Q_{N_{\mathrm{cov}}}(v)=s_{N_{\mathrm{cov}}},
\qquad
Q_j(v)=
\begin{cases}
Q_{j+1}(v),&s_j=s_{j+1},\\[2mm]
\displaystyle
\left(v-\frac{c_j+c_{j+1}}2\right)Q_{j+1}(v),
&s_j\ne s_{j+1},
\end{cases}
\quad(1\le j<N_{\mathrm{cov}}),
\qquad
Q_{\mathrm{sgn}}=Q_1.
\]
\endgroup
Each sign change contributes one linear factor, so
\[
\deg Q_{\mathrm{sgn}}\le d_{\mathrm{flip}}\le L.
\]
To check its sign, descending induction gives
\[
s_jQ_j(v)>0
\quad\text{if }v<c_j+\varepsilon_{\mathrm{ext}},
\]
and
\[
s_kQ_j(v)>0
\quad\text{if }j\le k\le N_{\mathrm{cov}}
\text{ and }|v-c_k|<\varepsilon_{\mathrm{ext}}.
\]
Indeed, the terminal polynomial is the constant \(s_{N_{\mathrm{cov}}}\). At an equal-sign pair the induction carries over unchanged. At an opposite-sign pair,
\[
c_{j+1}-c_j>\delta_{\mathrm{ext}}>4\varepsilon_{\mathrm{ext}}.
\]
The new linear factor is negative for \(v<c_j+\varepsilon_{\mathrm{ext}}\) and positive within \(\varepsilon_{\mathrm{ext}}\) of every later center. These signs, together with \(s_j=-s_{j+1}\), establish the induction.

For any \(v\in\mathcal K_{\mathrm{ext}}\), choose a covering center \(c_j\). Opposite signs at \(v\) and \(c_j\) would imply
\[
|v-c_j|>\delta_{\mathrm{ext}},
\]
contradicting \(|v-c_j|<\varepsilon_{\mathrm{ext}}\). Hence
\[
\rho(v)=s_j\mathcal E_{\mathrm{app}},
\qquad
\rho(v)Q_{\mathrm{sgn}}(v)
=
\mathcal E_{\mathrm{app}}\,s_jQ_{\mathrm{sgn}}(v)>0.
\]
This completes the sign-polynomial construction.
% lean: Causalean.Mathlib.Analysis.FinitePolynomialAlternationDuality.exists_signPolynomial_of_no_alternatingExtrema

\item We now turn that strict sign agreement into a uniform improvement. Define
\[
p_{\mathrm{sgn}}(v)=\rho(v)Q_{\mathrm{sgn}}(v),
\qquad
\mathcal B_{\mathrm{sgn}}
=
\{v\in[v_-,v_+]:p_{\mathrm{sgn}}(v)\le0\}.
\]
This set is compact and contains no point of \(\mathcal K_{\mathrm{ext}}\). Define
\[
c_{\mathrm{imp}}
=
\begin{cases}
\displaystyle
\frac{\max_{v\in\mathcal B_{\mathrm{sgn}}}|\rho(v)|
+\mathcal E_{\mathrm{app}}}{2},
&\mathcal B_{\mathrm{sgn}}\ne\varnothing,\\[3mm]
\mathcal E_{\mathrm{app}}/2,
&\mathcal B_{\mathrm{sgn}}=\varnothing.
\end{cases}
\]
In the first case the maximum is strictly below \(\mathcal E_{\mathrm{app}}\), since otherwise a maximizing point would belong to \(\mathcal K_{\mathrm{ext}}\). Thus in either case
\[
c_{\mathrm{imp}}<\mathcal E_{\mathrm{app}},
\qquad
|\rho(v)|\ge c_{\mathrm{imp}}
\ \Longrightarrow\
p_{\mathrm{sgn}}(v)>0
\quad(v\in[v_-,v_+]).
\]
Define
\begingroup\small
\[
\mathcal H_{\mathrm{imp}}
=
\{v\in[v_-,v_+]:|\rho(v)|\ge c_{\mathrm{imp}}\},
\qquad
m_{\mathrm{imp}}
=
\min_{v\in\mathcal H_{\mathrm{imp}}}p_{\mathrm{sgn}}(v),
\qquad
M_{\mathrm{imp}}
=
\max_{v\in[v_-,v_+]}|Q_{\mathrm{sgn}}(v)|.
\]
\endgroup
The set \(\mathcal H_{\mathrm{imp}}\) is compact and contains \(\mathcal K_{\mathrm{ext}}\). Therefore \(m_{\mathrm{imp}}>0\). Also \(M_{\mathrm{imp}}>0\), because \(Q_{\mathrm{sgn}}\) is nonzero at every extremum. Set
\[
t_{\mathrm{imp}}
=
\min\left\{
\frac{m_{\mathrm{imp}}}{M_{\mathrm{imp}}^2},
\frac{\mathcal E_{\mathrm{app}}-c_{\mathrm{imp}}}{2M_{\mathrm{imp}}}
\right\}>0.
\]
For \(v\in\mathcal H_{\mathrm{imp}}\),
\[
t_{\mathrm{imp}}Q_{\mathrm{sgn}}(v)^2
\le t_{\mathrm{imp}}M_{\mathrm{imp}}^2
\le m_{\mathrm{imp}}
\le p_{\mathrm{sgn}}(v)
<2p_{\mathrm{sgn}}(v).
\]
Consequently,
\[
\begin{aligned}
\bigl(\rho(v)-t_{\mathrm{imp}}Q_{\mathrm{sgn}}(v)\bigr)^2
&=\rho(v)^2-2t_{\mathrm{imp}}p_{\mathrm{sgn}}(v)
+t_{\mathrm{imp}}^2Q_{\mathrm{sgn}}(v)^2\\
&<\rho(v)^2
\le\mathcal E_{\mathrm{app}}^2.
\end{aligned}
\]
For \(v\in[v_-,v_+]\setminus\mathcal H_{\mathrm{imp}}\),
\[
\begin{aligned}
|\rho(v)-t_{\mathrm{imp}}Q_{\mathrm{sgn}}(v)|
&\le|\rho(v)|+t_{\mathrm{imp}}|Q_{\mathrm{sgn}}(v)|\\
&<c_{\mathrm{imp}}
+\frac{\mathcal E_{\mathrm{app}}-c_{\mathrm{imp}}}{2}
<\mathcal E_{\mathrm{app}}.
\end{aligned}
\]
The perturbed residual is continuous on the compact interval, so these pointwise strict inequalities imply
\[
\|\rho-t_{\mathrm{imp}}Q_{\mathrm{sgn}}\|_\infty
<\mathcal E_{\mathrm{app}}.
\]
But the polynomial
\[
Q_{\mathrm{imp}}=Q_*+t_{\mathrm{imp}}Q_{\mathrm{sgn}}
\]
has degree at most \(L\), and hence
\[
\mathcal E_{\mathrm{app}}
\le\|g-Q_{\mathrm{imp}}\|_\infty
=\|\rho-t_{\mathrm{imp}}Q_{\mathrm{sgn}}\|_\infty
<\mathcal E_{\mathrm{app}},
\]
a contradiction.

Thus, when \(\mathcal E_{\mathrm{app}}>0\), there exist increasing interval nodes \(v_1,\ldots,v_{L+2}\) with alternating residual values. Define their orientation by
\[
\eta=\frac{\rho(v_1)}{\mathcal E_{\mathrm{app}}}\in\{-1,1\}.
\]
Together with the zero-error construction, we have in both cases
\[
v_-\le v_1<\cdots<v_{L+2}\le v_+,
\qquad
\eta\in\{-1,1\},
\qquad
\rho(v_i)=\eta(-1)^{i-1}\mathcal E_{\mathrm{app}}.
\]
% lean: Causalean.Mathlib.Analysis.FinitePolynomialAlternationDuality.exists_strict_uniformImprovement
% lean: Causalean.Mathlib.Analysis.FinitePolynomialAlternationDuality.exists_equioscillationWitness

\item Define the nodal weights, their normalizing sum, and the normalized weights by
\[
\beta_i
=
\prod_{\substack{1\le j\le L+2\\j\ne i}}(v_i-v_j)^{-1},
\qquad
Z_{\mathrm{norm}}=\sum_{i=1}^{L+2}|\beta_i|,
\qquad
\widetilde w_i=\frac{\beta_i}{Z_{\mathrm{norm}}}
\quad(1\le i\le L+2).
\]
Distinctness of the nodes makes every factor nonzero, so \(Z_{\mathrm{norm}}>0\). Hence
\[
\sum_{i=1}^{L+2}|\widetilde w_i|
=
\frac{\sum_{i=1}^{L+2}|\beta_i|}{Z_{\mathrm{norm}}}
=1.
\]
There are exactly \(L+2-i\) negative factors in \(\beta_i\). Therefore
\[
\beta_i=(-1)^{L+2-i}|\beta_i|,
\qquad
\widetilde w_i=(-1)^{L+2-i}|\widetilde w_i|.
\]
% lean: Causalean.Mathlib.Analysis.FinitePolynomialAlternationDuality.sum_abs_normalizedLagrangeWeight_eq_one
% lean: Causalean.Mathlib.Analysis.FinitePolynomialAlternationDuality.normalizedLagrangeWeight_alternates

\item For every real polynomial \(Q\) of degree at most \(L\), Lagrange interpolation at these \(L+2\) nodes gives
\[
Q(v)
=
\sum_{i=1}^{L+2}
Q(v_i)
\prod_{\substack{1\le j\le L+2\\j\ne i}}
\frac{v-v_j}{v_i-v_j}
\quad(v\in\mathbb R).
\]
The coefficient of \(v^{L+1}\) on the left is zero, whereas on the right it is \(\sum_i\beta_iQ(v_i)\). Dividing by \(Z_{\mathrm{norm}}\) yields
\[
\sum_{i=1}^{L+2}\widetilde w_iQ(v_i)=0.
\]
In particular,
\[
\sum_{i=1}^{L+2}\widetilde w_i v_i^h=0
\quad(0\le h\le L),
\qquad
\sum_{i=1}^{L+2}\widetilde w_iQ_*(v_i)=0.
\]
% lean: Causalean.Mathlib.Analysis.FinitePolynomialAlternationDuality.sum_normalizedLagrangeWeight_mul_pow_eq_zero
% lean: Causalean.Mathlib.Analysis.FinitePolynomialAlternationDuality.sum_normalizedLagrangeWeight_mul_eval_eq_zero

\item Define the weighted target value by
\[
a_{\mathrm{dual}}
=
\sum_{i=1}^{L+2}\widetilde w_i g(v_i).
\]
Polynomial cancellation and the alternating residual identity give
\[
\begin{aligned}
a_{\mathrm{dual}}
&=
\sum_{i=1}^{L+2}\widetilde w_i\bigl(g(v_i)-Q_*(v_i)\bigr)
+\sum_{i=1}^{L+2}\widetilde w_iQ_*(v_i)\\
&=
\sum_{i=1}^{L+2}
(-1)^{L+2-i}|\widetilde w_i|\,
\eta(-1)^{i-1}\mathcal E_{\mathrm{app}}\\
&=
\eta(-1)^{L+1}\mathcal E_{\mathrm{app}}
\sum_{i=1}^{L+2}|\widetilde w_i|\\
&=\eta(-1)^{L+1}\mathcal E_{\mathrm{app}}.
\end{aligned}
\]
Since \(\eta\in\{-1,1\}\) and \(\mathcal E_{\mathrm{app}}\ge0\), this proves
\[
|a_{\mathrm{dual}}|=\mathcal E_{\mathrm{app}}.
\]
% lean: Causalean.Mathlib.Analysis.FinitePolynomialAlternationDuality.exists_alternationDualCertificate
% lean: Causalean.Mathlib.Analysis.FinitePolynomialAlternationDuality.AlternationDualCertificate.toFiniteMomentDual
% lean: Causalean.Mathlib.Analysis.FinitePolynomialAlternationDuality.exists_finiteMomentDual

\item Finally, choose the orientation of the weights by
\[
w_i=
\begin{cases}
\widetilde w_i,&a_{\mathrm{dual}}\ge0,\\
-\widetilde w_i,&a_{\mathrm{dual}}<0,
\end{cases}
\quad(1\le i\le L+2).
\]
If \(a_{\mathrm{dual}}\ge0\), the preceding identities directly give
\[
\sum_{i=1}^{L+2}|w_i|=1,
\qquad
\sum_{i=1}^{L+2}w_iv_i^h=0
\quad(0\le h\le L),
\qquad
\sum_{i=1}^{L+2}w_ig(v_i)
=a_{\mathrm{dual}}
=|a_{\mathrm{dual}}|
=\mathcal E_{\mathrm{app}}.
\]
If \(a_{\mathrm{dual}}<0\), negating every weight preserves its absolute value and reverses every weighted sum:
\[
\sum_{i=1}^{L+2}|w_i|
=\sum_{i=1}^{L+2}|-\widetilde w_i|=1,
\qquad
\sum_{i=1}^{L+2}w_iv_i^h
=-\sum_{i=1}^{L+2}\widetilde w_iv_i^h=0,
\]
and
\[
\sum_{i=1}^{L+2}w_ig(v_i)
=-a_{\mathrm{dual}}
=|a_{\mathrm{dual}}|
=\mathcal E_{\mathrm{app}}.
\]
Thus both cases establish the required weights, including the case of zero approximation error.

For the signed-measure formulation, define
\[
\sigma=\sum_{i=1}^{L+2}w_i\delta_{v_i}.
\]
Because the nodes are distinct, for every real-valued function \(f\) on the interval,
\[
\int f(v)\,d\sigma(v)=\sum_{i=1}^{L+2}w_if(v_i),
\qquad
\|\sigma\|_{\mathrm{TV}}=\sum_{i=1}^{L+2}|w_i|=1,
\qquad
\operatorname{supp}\sigma\subseteq\{v_1,\ldots,v_{L+2}\}.
\]
Applying the integral identity to \(f(v)=v^h\) and to \(f=g\) gives the asserted moments and target value.
% lean: CausalSmith.Stat.AnnotationRarearmFrontier.FinitePolynomialDuality
\end{enumerate}
\end{proof}
